\section{Metodología}
Este trabajo se base en aplicar el método experimental-desarrollador que se enfoque en crear, diseñar y modelar de algoritmos inteligentes y luego hacer las pruebas necesarias para validar y evaluar su eficacia en un entorno que simula condiciones reales, este trabajo se describe de la siguiente manera:
\begin{itemize}
  \item Primera fase la Recolección y análisis de datos: Identificar de las fuentes de datos que vamos a usar conjuntos de datos abiertos confiables. Y haremos el procesamiento de datos limpiar los datos y transformarlos en formatos analizables para descubrir los patrones históricos del consumo de los recursos.
  \item Segunda fase el diseño del modelo de predicción: elegir algoritmos adecuados que vamos a crear un modelo que se base en aprendizaje profunda o las redes neuronales por su potencia elevada de predecir de series de tiempo para (workload pattrens). usaremos los datos históricos en el entrenamiento del modelo para predecir (peak times) 'periodos de inactividad' (idle times) en la demanda de datos.
  \item Tercera fase la simulación y los entornos: por la dificultad ir presencialmente a los centros de datos, usaremos herramientas de simulación como Cloudsim o iFogSim para construir entornos de nube virtual así podremos aplicar algoritmos para las gestiones de recursos que encenderán y apagarán los servidores virtuales (VMs) según al resultado de los modelos generados anteriormente.
\item Ultima fase la evaluación y validación: en esta fase evaluaremos y validaremos Energy Saving, Prediction Accuracy, SLA Violation.
\end{itemize}
